%==============================================================================
\section{Routing in depth}\label{sec:routing}
%==============================================================================

Routing, Eq.~\eqref{eq:psi}, is the step that takes most of the design effort in the code: it is a graph problem, not
an array operation, its cost and robustness depend on the traversal, and real ice-sheet data make the graph cyclic. This
section explains the four traversals, why the default is to record and replay, what cycles do, and how conservation
is guaranteed.

\subsection{The routing graph}\label{rt:graph}
Draw an edge $n\to c$ between neighbouring grounded cells when $w_{n\to c} > 0$. Then $\psiout(c)$ depends on
$\psiout(n)$ for every edge $n\to c$. If the graph is a DAG (no cycle), Eq.~\eqref{eq:psi} has a unique solution, obtained
by evaluating cells in any topological order. If it has a cycle, Eq.~\eqref{eq:psi} is a linear system around the cycle,
which none of the traversals solves; they each produce some approximation.

\subsection{Four traversals}\label{rt:algos}
\begin{center}\small
\begin{tabular}{>{\raggedright\arraybackslash}p{2.6cm}>{\raggedright\arraybackslash}p{5.2cm}>{\raggedright\arraybackslash}p{7.2cm}}
\toprule
algorithm & how & properties \\
\midrule
\code{RecursivePsiOut} & depth-first recursion from each grounded cell, \code{visited} flags (\code{accumulate\_psi\_out!}) & KORI-ULB's
  algorithm. Recurses as deep as the longest flow path: on real grids this can exceed the process stack (\code{ulimit -s}) and crash
  with a \code{Stack\allowbreak{}Overflow\allowbreak{}Error}, before the visit cap binds. GDSWarner only. \\
\code{IterativePsiOut} & the same traversal with an explicit stack (\code{update\_psi\_out\_iterative!}) & Same field bit for bit,
  no depth limit; the source measured it 30--40 times slower per sweep than the recursion on Thwaites 2\,km and pan-Antarctica
  8\,km. GDSWarner only. \\
\code{Topological\allowbreak{}Psi\allowbreak{}Out} & Kahn's algorithm on the edge list (\code{update\_psi\_out\_topological!}) & $O(n)$, no cap, no stack. Exact only on a
  DAG; errors on a cycle unless \code{allow\_cycles = true}. GDSWarner only. \\
\code{TapedPsiOut} (default) & record the recursion's operations once per \code{update\_q!}, replay them every sweep & Same field as the recursion
  bit for bit, including partial reads in cycles and the cap; no depth limit; every scheme. \\
\bottomrule
\end{tabular}
\end{center}
\paragraph{Why record and replay.} The order in which the recursion visits cells and folds in contributions depends only on
the mask and the fractions $w$, i.e.\ on the geometry, never on the source. Within one \code{update\_q!} the geometry is fixed, but
the dissipation and $(q,N)$ loops route the same graph with a new $\mdot_{\rm total}$ up to 20 times. \code{TapedPsiOut} records the
traversal as a flat list of \code{dst += src * w} and \code{dst = max(0, dst)} operations (with an explicit stack, so no depth
limit) and each sweep replays it: a single linear pass over five vectors, with no branching on the graph. The tape is
invalidated at the start of every \code{update\_q!}, since the gradients are recomputed there. Recording applies every source up
front rather than on first visit; this changes no value read, because a cell is never read before its first visit. A
replay is linear in the source as long as no clamp binds (the tour doubles the source and gets exactly double
$\psiout$).

\paragraph{The explicit stack.} A frame is $(i, j, k)$: $k = 0$ means not yet visited; $1\le k\le K$ means neighbour $k$ is next
($K = 4$, or 8 for 8-neighbour schemes); $k = K+1$ means done (clamp, pop). When neighbour $k$ is unvisited, it is pushed
\emph{without} advancing the parent's $k$; after it is resolved and popped, the parent finds the same neighbour, now visited,
and folds its value in. This is exactly how a return value reaches a paused caller in the recursion, which is why the
results agree bit for bit.

\subsection{Cycles}\label{rt:cycles}
In the recursion, a cell on the stack whose descendant routes back to it is read while still incomplete: the descendant
receives a partial value, with no warning. The recursion never fails on a cycle, but the result depends on the visiting
order. A topological sort cannot process the cells of a cycle at all (their in-degree never reaches zero), which is why
\code{TopologicalPsiOut} errors by default: the docstring records a case where silently using the truncated field made the
correlation of $q$ with the recursion's fall to near zero.

Which graphs have cycles:
\begin{itemize}[leftmargin=1.5em]
\item GDS schemes: the fractions come from the smoothed gradient components, normalised per cell, which is not the gradient of
  any one surface. The source records cycles on roughly half the grounded cells of Thwaites 2\,km at \code{longcoupwater} = 5,
  more (37\,514 of 51\,945) with no smoothing at all, and 2573 of 3242 on an 8\,km coarsening: real, noisy topography alone is
  enough. Smoothness alone does not prevent them: the two cells either side of a valley line each have a centred
  gradient pointing at the other, a 2-cycle. The tour's bump-free fixture, whose thickness varies as $\sin(y)$, has
  such a valley and the topological sort refuses it; with a uniform flow direction the graph is acyclic and the tour
  checks topological = recursive there.
\item Warner, Quinn and Tarboton: water goes only to strictly lower cells of $\phif$ (Tarboton's chosen facet has both slopes
  positive, or sends everything to its one lower neighbour), so there is no cycle; the test suite checks conservation for all
  three. ModifiedTarboton takes its direction from the centred gradient of $\phif$ and can send water to a higher neighbour.
\end{itemize}
The water lost to cycles is the main reason Warner replaced GDSWarner as the default: the source gives 3\,\%, 8\,\% and 34\,\% of
the positive basal melt lost with GDSWarner on Thwaites 2\,km, Greenland 8\,km and Antarctica 16\,km, and none with Warner.

\subsection{Conservation}\label{rt:conservation}
With Warner (or Quinn, or Tarboton) and \code{PriorityFloodFill}, routing conserves water exactly: every grounded cell has a strictly lower
neighbour or is an outlet, the graph is a DAG, and each cell's outflow fractions sum to 1 or 0. Summing Eq.~\eqref{eq:psi}
over the grounded cells, every routed volume ends either in a non-grounded cell (across the grounding line or ice margin),
or at a grounded edge cell with nowhere lower to go (a sink, which keeps it). The tour checks
\[
\sum_{c\ \rm grounded} \frac{\mdot_c\Dx\Dy}{\rhow}
= \sum_{c} \psiout(c)\sum_{d:\ \text{off-domain or not grounded}} w_{c,d} \;+\; \sum_{c:\ \sum_d w_{c,d} = 0}\psiout(c)
\]
to $10^{-12}$ on its bumpy fixture, provided no clamp binds ($\mdot\ge0$ everywhere) and the visit cap is not reached. With
face fluxes the same balance holds cell by cell: the net outflow through a cell's four faces equals its source.

Two things break conservation by design: the clamp $\psiout\ge0$ (a refreezing cell destroys the inflow it refreezes; this
is physical, and \code{freeze\_on\_capacity!} reports the water available for it), and the visit cap.

\subsection{Cost}\label{rt:cost}
Recording costs one traversal with stack manipulation and vector pushes; replaying is a tight loop over the operation list,
whose length is about the number of grounded-to-grounded edges with $w>0$ plus one clamp per cell. For the timings of the
whole solve on the real grids, see Section~\ref{life:timing}.
